\section{Collection and document management: organisation, classification and
lifecycle}
\label{sec:organization}

Organisation determines which collections exist, which tier can see a
document, and whether the stored state remains trustworthy. HMRS places
documents by content rather than filename; split parts inherit the source
assignment, and the catalogue exposes library and document descriptions
before full bodies are read. This section covers the collection and document
levels of Table~\ref{tab:levels}; the segment level is treated in
Section~\ref{sec:retrieval}.

\subsection{Classification, metadata, tags and duplicate review}
\label{sec:org-placement}
\label{sec:org-registries}
\label{sec:org-tags}
\label{sec:org-dedup}

Integrity checks compare three sources: the document files on disk, the
\texttt{.tree-fs.json} tree and each library's \texttt{.knowledge-base.yml}.
The comparison covers document existence, path and index metadata. Tags derive
from content and are checked against it.
Duplicate detection reports candidate pairs for review; deletion is not
automatic because adjacent parts may be similar while containing distinct
source sections.

The layer's operations are library and sub-library creation, classification,
description and tag edits, move, rename, delete, duplicate review and forced
re-index. Together they implement the management loop that turns parsed content
into a maintained collection: placement records ownership, descriptions expose
scope, lifecycle operations change ownership or content, and reconciliation
checks whether the representations still agree. Each operation updates the
state relevant to it; for example, a move changes the document path and
registry entries, while re-indexing updates search indices. The audit compares
the files and registries where applicable. The catalogue lists libraries, their
descriptions, sub-library containers and document descriptions, which lets a
reader or an agent decide where to look before opening a body. Classification
involves a semantic choice: the ingest agent
selects the library from the parsed content, and the platform checks the
consequences rather than the choice. The same asymmetry holds at query time: the workflow selects a collection,
whereas operations to open and index documents, and optionally score their
segments, are provided by the server.

Sub-libraries are tree nodes with their own descriptions and indices. For a
parent scope, the vector service resolves the parent and its descendants to the
collections available to that query; the two-stage service likewise resolves
descendant identifiers during its lexical stage. This is a scope-resolution
policy, not a claim that child vectors are stored in the parent's collection.
The measured campaigns used explicit per-library scopes for the flat union and
did not treat parent traversal as an independently benchmarked retrieval
condition. Index availability and route-specific failures therefore still
require a post-search probe.

\subsection{Lifecycle operations, graph and operational knowledge}
\label{sec:org-lifecycle}
\label{sec:org-graph}
\label{sec:org-experience}

A historical rename trial left stale vector chunks until forced re-indexing.
The current rename path starts index synchronisation asynchronously, with
failures logged as non-fatal; checking the resulting index remains necessary
because successful metadata updates alone do not establish index consistency. The document graph supports navigation and candidate
expansion, not evidence. Operational lessons are stored as retrievable
documents with credibility tiers, but the available measurements are case
examples rather than a scale evaluation of distillation, approval or
credibility decay.

\subsection{Organisation contract}
\label{sec:org-integrity}
\label{sec:org-summary}

The evaluation checks three-source consistency, content-derived descriptions,
regression probes and lifecycle outcomes (Section~\ref{sec:e2}). Reconciliation
can expose selected stale or inconsistent states; it does not detect or prevent
every failure by construction. Merge and split remain outside the controlled test
scope, so performance claims must be read within the corpus sizes reported in
the Evaluation section.
